%% =====================================================================================
%% 2 Related work. Each paragraph ends with how this paper differs.
%% The last paragraph declares the overlap with a concurrent manuscript. It is required
%% by the dual-submission policy; do not cut it to save space. Cite it in the third
%% person only (double-blind).
%% Claims about other papers were checked against their full text on 2026-09-16.
%% =====================================================================================
\section{Related work}
\label{sec:related}

\paragraph{Collective risk and selection.}
The collective-risk dilemma was built to study why groups fail to avoid a shared
disaster~\cite{milinski2008collective}. Human groups reach the target mostly when the risk
is high, and inequality and communication change the result~\cite{tavoni2011inequality}.
Games with a threshold have many equilibria~\cite{palfrey1984participation}, human
contributions decline over the rounds of repeated public goods games~\cite{andreoni1988free},
and the risk of collective failure can turn a defection trap into a coordination
problem~\cite{santos2011risk}. Theory also shows that selection on payoff relative to a small
group favours whoever does better than its neighbours~\cite{schaffer1988evolutionarily},
while selection between groups can favour cooperation~\cite{traulsen2006evolution}. This
literature tells us that behaviour \emph{should} move with the risk, so a flat response is
informative, and we reuse its selection theory to read our mixed tables.

\paragraph{Language models in games.}
Language models have been studied as economic
subjects~\cite{horton2023large,chen2023emergence,mei2024turing}, in repeated two-player
games~\cite{akata2025playing,fontana2024nicer}, and in multi-player game
benchmarks~\cite{huang2025competing}. Models tend to be more cooperative than
people~\cite{mei2024turing}, an even split can reflect a focal point rather than
fairness~\cite{dodivers2025uncovering}, and models can hold a correct belief or plan and act
against it~\cite{fan2024can,schmied2026llms}. In social dilemmas, societies of agents
over-use a common resource~\cite{piatti2024cooperate}, reasoning models free-ride in public
goods games~\cite{guzmanpiedrahita2025corrupted}, and agents replace people in a
collective-risk game~\cite{slumbers2025using}. Closest to us, Willis et al.\ study the same
collective-risk game with strategies written by three models, map welfare over mixtures of
selfish and collective strategies, and find that payoff-biased imitation drives larger
groups to selfish, low-welfare outcomes~\cite{willis2026evaluating}. We differ in three ways.
We record the turn-by-turn play of deployed models, not code they write. We price that play
against a computed optimum, with two tests where paying cannot help. And we show that
whether selection compares tablemates or separate groups decides which model wins, which
gives one simple route to selfish outcomes.

\paragraph{Multiagent evaluation.}
Sequential social dilemmas made cooperation a property of learned
policies~\cite{leibo2017multi,hughes2018inequity}, and cooperative AI asks for agents that
work with partners unlike themselves~\cite{dafoe2020open,stone2010ad}. Melting Pot tests a
focal population against background agents it has never met~\cite{leibo2021scalable}.
Empirical game-theoretic analysis treats trained agents as strategies of a
meta-game~\cite{wellman2006methods,tuyls2018generalised}, and $\alpha$-Rank ranks them by
evolutionary dynamics~\cite{omidshafiei2019alpha}. Tournaments that mix language models with
classic strategies report model-specific strategic fingerprints in the prisoner's
dilemma~\cite{payne2025strategic}, and cooperation evolves differently across base models in
a donor game~\cite{vallinder2025cultural}. We add a many-player dilemma in which the best
response to fixed partners can be computed exactly, and real tables of two different models.

\paragraph{Overlap with concurrent work.}
A concurrent manuscript uses the same game and an overlapping set of models to ask whether
models reproduce the risk sensitivity of human players, and suggests that more capable models
can respond to risk~\cite{anon-if}. It reports a self-play risk baseline related to part of
Section~\ref{sec:risk}. The present paper asks how low-cost agents behave against the
optimum, against fixed partners and against other models; the tests where paying cannot
help and Sections~\ref{sec:knowing} to~\ref{sec:selection} have no counterpart there.
